The default definition is part of the PD, and the treatment of withdrawn ratings and the size of the cohort are part of it too. On the corporate cohorts of S\&P from 2000 to 2025 the long-run PD of CCC to C is $\spLraDtwoCcc\%$ on the default page and $\spLraDfourCcc\%$ on the transition page (gap index $\spPtRatioCcc$, 95\% interval $\spRatioCccLo$ to $\spRatioCccHi$), and the risk weight of the cohort mix falls from $\spRwDtwo\%$ to $\spRwDfour\%$, a fall of $\spRwGap$ points with a 95\% interval of $\spRwGapLo$ to $\spRwGapHi$. On rating paths of known truth with withdrawal ten times as likely in the year before default, removing the withdrawn ratings understates the PD of grade B by $\removedBkNine\pm\removedBkNinese$ percent, keeping them by $\keptBkNine\pm\keptBkNinese$ percent, following them to the year end by $\followedBkNine\pm\followedBkNinese$ percent and treating them as censored in a duration estimator by $\durBkNine\pm\durBkNinese$ percent, and the risk weight of the cohort mix is $\rwAbsremovedkNine$ points too low under removal; the tests of time homogeneity and of the last matrix reject $\testthkNinefive\%$ and $\testmrkNinefive\%$ of the repetitions at $k=9$, and the second rejects $\testmrkNineteenfive\%$ at $k=19$, so that a matrix test sees the withdrawals, but late. On the years in which Fitch's transition page has a populated default category the risk weight of its cohort mix falls by $\fiPopRwGap$ points between the pages, against $\fiRwGap$ points on all years. With 50 obligors per grade the Wald interval contains the PD of AAA in $\waldAaaFifty\%$ of the samples, and the enumerated coverage of the other intervals is nominal only at the cost of widths of hundreds of basis points.

The theory explains the pattern. Removing or keeping a withdrawn rating is biased on the first cohort by a factor that depends on $w(1+k)$ and, to first order in the PD, not on the PD (the first-order formula is refuted for the grade CCC to C, where the exact closed form is needed); following it is exact on the first cohort and biased on every later one, because the cohort of the next year is already short of the ratings about to default; and the coverage of a Wald interval for a rare PD is at most the probability of observing a default. The closed forms of the first cohort agree with the simulation within $\thMaxZFirst$ standard errors over $\thNcmpFirst$ distinct comparisons, and the numerical evaluation of the pooled expectations within $\thMaxZPooled$ over $\thNcmpPooled$; for a multi-state grade the two-state numerical evaluation misses the followed treatment, in opposite directions for B and for CCC to C.

The next steps are three. The protocol should be run on a bank's own default and withdrawal history, where the follow-up that the guidelines ask for may exist and the informativeness can be estimated rather than assumed. An estimator that anchors the informativeness on a second count of defaults, or reports the PD as a function of it, should be evaluated against the biases measured here. And the bootstrap of the real-data gap should be repeated with a block resample of the cohort years, to state how much of its width the serial dependence of the cycle would add.
